\chapter{Conclusions and Recommendations}

\section{Conclusions}

This work presents Omnisoot as a robust computational tool built on Cantera to simulate soot formation coupled with gas-phase chemistry across various combustion and pyrolysis scenarios. Omnisoot integrates four reactor models, two particle dynamics models, and four inception and surface growth models, providing a comprehensive platform for analyzing uncertainties in predicting soot yield and morphology originating from chemical kinetics, inception processes, and surface growth rates.

Omnisoot was employed to elucidate differences among inception models in shock tubes, flow reactors, and perfectly stirred reactors. Global adjustment factors were applied to inception and surface growth rates to accurately reproduce experimental data for each target scenario across a range of parameters, including temperature, pressure, composition, and flow rates. The performance of inception models was evaluated based on soot yield, morphology, size distribution, inception and surface growth rates, and carbon mass source analysis. Although different inception models predicted similar carbon yields that closely matched experimental observations, they produced varying morphological characteristics. This highlights the importance of characterizing soot morphology to better constrain inception flux and surface growth rate predictions.

Simulations of ethylene pyrolysis in a flow reactor demonstrated the temperature sensitivity of inception models. Irreversible models (Irreversible Dimerization and Dimer Coalescence) effectively captured the bimodal PSD resulting from high inception flux in the lower-temperature region near the reactor exit due to their weak dependence. However, the Reactive Dimerization and E-Bridge models generated nearly unimodal distributions due to their exponential temperature dependence, causing a significant drop in inception flux in the same region.

The PSD sampled at the outlet of a flow reactor downstream of a perfectly stirred reactor showed good agreement with experimental measurements across three equivalence ratios. Simulations indicated that active inception and surface growth primarily occur at high temperatures within the perfectly stirred reactor and at the entrance region of the flow reactor, far upstream of the sampling location. Under these conditions, the particle size distribution, morphology, and inception flux predictions were largely consistent across all inception models, although notable differences were observed in soot volume fraction due to variations in predicted HACA growth rate.


The imposed reactor model in Omnisoot was used to simulate soot formation during shock-tube pyrolysis of 5\%, 10\%, and 30\% methane over a temperature range of 1800–2500~K and pressures near 4~atm. The simulation results were compared with time-resolved measurements of key gas-phase species and soot volume. The performance of several kinetic mechanisms were evaluated, indicating that FFCM2 and CRECK reproduced the time history of CH$_4$, C$_2$H$_4$, and C$_2$H$_2$ in close agreement with the measurements despite significant differences in pathways. The analysis suggests that methane conversion and intermediate yields strongly depend on temperature and fuel concentration, with significant model–measurement deviations for CH$_4$ mole fraction emerging at 30\% loading due to increased absorption interference. Carbon flux analysis highlighted the dominant role of C$_3$H$_3$ in channeling carbon toward C$_2$H$_2$ and PAHs, and the role of P-C$_3$H$_4$ as a key precursor of C$_3$H$_3$. CRECK coupled with E-Bridge Modified model captured the time and temperature trends of $f_v$ at 5\% and 10\% loadings, but a substantial overprediction with similarities in time trends is observed at 30\% loading, indicating potential missing chemistry or physical effects that could inhibit surface growth at higher loadings.


%The present work provides a combined numerical and experimental investigation of CB formation during shock-tube pyrolysis of 5\% CH$_4$ in argon at a post-reflected-shock pressure of P$_5$=4.5$\pm$0.5~atm and temperatures spanning 1800~K$<$T$_5$$<$2500~K. Multi-wavelength laser diagnostics were employed to measure the time-resolved mole fractions of CH$_4$, C$_2$H$_4$, and C$_2$H$_2$. The CB volume fraction, $f_v$, was determined using light extinction at 633~nm and 1064~nm, while collected TEM images were analyzed to extract the primary particle diameter, $d_p$, as a quantitative measure of CB morphology. 

Simulations using the CRECK mechanism accurately reproduced methane conversion and acetylene yields. The gas-phase chemistry  solver was coupled with a sectional population balance model that accounts for CB inception, surface growth, and the reduction in H/C due to carbonization.

Further analysis was performed for 5\% CH$_4$ pyrolysis by quantifying primary particle diameter using TEM images of samples collected after each shock-tube test. Simultaneous calibration of the model using both $f_v$ and $d_p$ provided strong constraints on inception, PAH adsorption, HACA surface growth, and carbonization rates, highlighting the importance of morphology data in avoiding non-unique calibrations based solely on $f_v$. The model reproduced the time history of $f_v$ and the final $d_p$ in good agreement with the measurements. Four stages of CB evolution were identified: (I) an induction period dominated by precursor formation, (II) rapid increase in $f_v$ driven by rising surface growth rates, (III) transition toward final yield, composition, and $d_p$ as surface growth wanes, and (IV) a coagulation-dominated stage in which primary particle diameters are chemically frozen while the agglomerate mobility diameter continues to increase.

The analysis showed that PAH adsorption dominates CB mass growth at lower temperatures and early times, whereas HACA becomes the primary contributor with increasing T$_5$. The model reproduced the bell-shaped temperature dependence of $f_v$ but overpredicted induction times at T$_5$$>$2200~K, suggesting the presence of missing high-temperature PAH formation pathways that would accelerate inception and PAH adsorption. TEM measurements revealed a decreasing trend in primary particle diameter with increasing T$_5$, while the model predicted a bell-shaped dependence.

The maturity index, $\beta$, inferred from the extinction ratio at 633~nm and 1064~nm, exhibited strong similarities to the predicted H/C ratio of CB particles, which is used in the model as a measure of surface reactivity. Both $\beta$ and H/C decrease rapidly at early times before approaching near-asymptotic values, with the rate of the initial decline increasing with T$_5$.

% Time-resolved measurements of key gas-phase species and soot volume fraction during shock-tube pyrolysis of 5\%, 10\%, and 30\% methane over a temperature range of 1800–2500~K and pressures near 4~atm. Using multi-wavelength laser diagnostics, we quantified the evolution of CH$_4$, C$_2$H$_4$, and C$_2$H$_2$. The results show that methane conversion and intermediate yields strongly depend on temperature and fuel concentration, with significant model–measurement deviations for CH$_4$ mole fraction emerging at 30\% loading due to increased absorption interference. Comparisons with detailed kinetic models revealed that the FFCM2 and CRECK models capture gas-phase species more accurately despite significant differences in pathways. Carbon flux analysis highlighted the dominant role of C$_3$H$_3$ in channeling carbon toward C$_2$H$_2$ and PAHs, and the role of P-C$_3$H$_4$ as a key precursor of C$_3$H$_3$. The model captured the time and temperature trends of $f_v$ at 5\% and 10\% loadings, but a substantial overprediction with similarities in time trends is observed at 30\% loading, indicating potential missing chemistry or physical effects that could inhibit surface growth at higher loadings.

\section{Recommendations}

\subsection{Implementation of reactor network}
The current structure of Omnisoot provides four reactor models, including two close (batch) and two open (continuous) reactors. While single reactor are useful tools for description of lab-scale reactive systems under controlled conditions, large-scale industrial devices with complex flow fields are often represented as a series of interconnected open reactor models~\citep{savarese2024model}. Similar levels of complexity are also encountered in plasma pyrolysis processes used for CB synthesis~\citep{lawson2025assessing}. Therefore, the development and implementation of a reactor network framework within Omnisoot is an essential next step to enable more realistic process representation and to support reactor design and optimization studies.

\subsection{Extending the oxidation model}

The oxidation model in Omnisoot describes the removal of carbon atoms through surface reactions with OH and O$_2$. However, it does not account for diffusion of O$_2$ into the soot bulk and the resulting internal oxidation, which becomes significant at temperatures between 600 and 1100~K. There is also growing evidence that soot gasification by CO$_2$ can reduce primary particle size and alter the nanostructure of soot~\citep{chang2020experimental}. Although these processes do not directly affect the current predictive capability of Omnisoot for CB synthesis conditions, extending the oxidation model to include such effects is an important step toward ensuring its applicability across a wider temperature range and for different carrier gas compositions.

\subsection{Developing the flame solver}

Soot formation in laminar burner-stabilized premixed and counterflow diffusion flames has been extensively studied in the literature. Standardized datasets are available on the ISF website\footnote{\href{https://www.adelaide.edu.au/cet/isfworkshop/data-sets}{https://www.adelaide.edu.au/cet/isfworkshop/data-sets}}, which provide a wide range of measurements, including species mole fractions, particle size distributions, soot volume fraction, and primary particle diameter. Accurate simulation of these flame configurations requires a one-dimensional solver, which is not currently available in Omnisoot. Development of such a solver is underway and is presently in the preliminary stages of implementation and validation.